% ============================================================
% ΛCDM+S: CWSF Written Report (science-fair version)
% Paste into Overleaf and compile with pdfLaTeX.
% ============================================================
\documentclass[12pt,letterpaper]{article}

\usepackage[margin=1in]{geometry}
\usepackage[T1]{fontenc}
\usepackage[utf8]{inputenc}
\usepackage{lmodern}
\usepackage{amsmath,amssymb}
\usepackage{booktabs}
\usepackage{graphicx}
\usepackage{setspace}
\usepackage{microtype}
\usepackage{xcolor}
\usepackage{caption}
\usepackage{enumitem}
\usepackage{float}
\usepackage{hyperref}

\setstretch{1.25}
\setlength{\parskip}{0.28em}
\setlength{\parindent}{1.5em}

\hypersetup{
  colorlinks=true,
  linkcolor=blue!50!black,
  citecolor=blue!50!black,
  urlcolor=blue!50!black
}

\newcommand{\Lcdms}{$\Lambda$CDM+S}
\newcommand{\tcrit}{t_{\mathrm{crit}}}

\title{%
  \textbf{\Lcdms: Does Extra Physics Help Fit the Universe?}\\[0.4em]
  \large A Computational Test of a Four-Parameter Cosmology Model\\
  Using Bayesian Inference
}
\author{Liam Desre}
\date{2026}

\begin{document}
\maketitle

\begin{abstract}
The standard model of cosmology, $\Lambda$CDM, describes a universe that
contains matter, radiation, and a constant dark-energy term $\Lambda$.
This project tests a closely related model called \Lcdms. The extra
``S'' is a smooth switch that turns the dark-energy term on over cosmic
time, instead of assuming it was always fully present. The physics of
that switch is locked in a separate theory file. This report is about
the computing test: I built a pipeline that predicts distances and
expansion rates, compares those predictions with supernovae, BAO, and
compressed CMB data, and then asks whether the two extra parameters are
worth the extra complexity.

I sampled four numbers---today's expansion rate $H_0$, today's dark-energy
fraction $\Omega_\Lambda$, a switch speed $k$, and a midpoint time
$\tcrit$---with Markov chain Monte Carlo. Nested sampling estimated the
Bayesian evidence. A recovery test showed that when the switch is turned
off, the code reproduces flat $\Lambda$CDM to better than one part in
$10^4$. On a published DES supernova comparison, \Lcdms\ fits the
points a little better ($\Delta\chi^2=-2.5$), but the usual complexity
penalties and the published evidence ratio still favour ordinary
$\Lambda$CDM. The computing stack therefore does what it is supposed to
do: it tests the model against data without claiming that a random walk
``proved'' the underlying physics.
\end{abstract}


%------------------------------------------------------------------------------
\section{Introduction}
%------------------------------------------------------------------------------

Most science-fair projects measure something in a lab. This one measures
a model. Cosmologists already have a successful recipe for the expanding
universe, called $\Lambda$CDM. It says that the universe is mostly dark
energy ($\Lambda$), plus cold dark matter and ordinary matter, and that
space is close to flat. That recipe fits a huge range of observations,
from nearby supernovae to the cosmic microwave background.

It is not perfect. Different experiments disagree about how fast the
universe is expanding today (the Hubble tension). There is also a
deeper conceptual question: why should $\Lambda$ be a fixed number that
was already sitting in the equations from the beginning? \Lcdms\ is a
proposed answer to the second question. In the theory layer of this
project, horizon entropy and a maximum-entropy endpoint are used to
argue that the expansion should settle toward a de~Sitter state---a
universe that expands at a constant rate. The computing layer does
\emph{not} re-derive that argument. It asks a more modest question:

\begin{quote}
If I allow dark energy to turn on smoothly, with two extra numbers
controlling the switch, do today's data prefer that story over ordinary
$\Lambda$CDM?
\end{quote}

That is a statistics problem, not a chalkboard-physics problem. The
honest way to answer it is to lock the physics, predict observables,
fit the data, and then penalize the extra freedom. This report
describes that pipeline in language a judge can follow, with only a
few equations.

%------------------------------------------------------------------------------
\section{Purpose and hypothesis}
%------------------------------------------------------------------------------

\paragraph{Purpose.}
Build and audit a complete inference stack for \Lcdms, then use it to
compare the model with $\Lambda$CDM on the same public datasets.

\paragraph{Hypothesis.}
If the extra switch is a real improvement, then after a fair
complexity penalty the data should prefer \Lcdms. If the extra
parameters only buy a slightly better curve-fit, information criteria
and Bayesian evidence should still prefer $\Lambda$CDM.

\paragraph{What this project is not.}
It is not a proof of Einstein's equation. It is not a claim that
``MCMC proved the attractor.'' Markov chain Monte Carlo only maps which
values of four parameters are allowed by the data. Nested sampling
estimates how much probability mass the whole model has, not whether
the second law of thermodynamics equals cosmic acceleration.

%------------------------------------------------------------------------------
\section{Background}
%------------------------------------------------------------------------------

\subsection{The standard picture in one page}

Imagine the universe as a balloon with dots on it. As the balloon
inflates, every pair of dots moves apart. The Hubble parameter $H$ is
the expansion rate. Redshift $z$ tells us how much light has been
stretched since it was emitted; larger $z$ means earlier times.
Supernovae give a distance-versus-redshift diagram. Baryon acoustic
oscillations (BAO) give a standard ruler in the galaxy distribution.
The cosmic microwave background (CMB) is leftover light from when the
universe was about 380{,}000 years old.

In flat $\Lambda$CDM, three ingredients set the expansion: radiation
(important early on), matter (which clumps), and a constant dark-energy
density (which dominates late). Today's matter fraction is then just
whatever is left after you choose today's dark-energy fraction.

\subsection{What ``+S'' adds}

\Lcdms\ keeps the same three ingredients, but multiplies the dark-energy
term by a fraction $\chi(t)$ that starts near zero and later approaches
one. Think of a dimmer switch, not an on/off light. The switch has a
speed $k$ and a midpoint time $\tcrit$ (the moment when the dimmer is
halfway up). When $k$ is sent to zero, the dimmer is stuck at fully on,
and the model must collapse back to ordinary $\Lambda$CDM. That recovery
is a computing test: if it fails, the code is not an extension of
$\Lambda$CDM and no posterior should be trusted.

The theory file also produces a thermodynamic equation of state for the
entropy sector. Computing treats that as a derived background array. It
is not a fifth sampled parameter.

\subsection{Bayesian inference without the jargon pile-up}

Three ingredients go into a Bayesian fit.

\begin{enumerate}[leftmargin=1.6em]
  \item \textbf{Prior.} What I believed about the four numbers
        \emph{before} looking at the likelihood. Here the priors are
        independent Gaussians (bell curves) whose widths come from
        mixing published $H_0$ information and from the switch timescale.
  \item \textbf{Likelihood.} How well a predicted catalogue matches the
        real catalogue, including measurement errors and, when
        available, covariance matrices.
  \item \textbf{Posterior.} The update: prior times likelihood. MCMC
        draws samples from that surface. Nested sampling integrates it
        to get the evidence $Z$, which is the right quantity for
        ``which model is better overall?''
\end{enumerate}

A better raw $\chi^2$ is not enough. Extra parameters can always hug
the points more tightly. AIC and BIC add a complexity fine. Evidence
goes further: it averages the likelihood over the whole prior, so a
model that only fits in a tiny corner of parameter space is
automatically down-weighted.

%------------------------------------------------------------------------------
\section{The model, in the few equations that matter}
%------------------------------------------------------------------------------

The computer only varies four numbers:
\begin{equation}
  \theta = (H_0,\; \Omega_\Lambda,\; k,\; \tcrit).
\end{equation}
Radiation today is held fixed. Matter today is whatever is left:
$\Omega_{m,0}=1-\Omega_\Lambda-\omega_{r,0}$. The dimmer has a closed
form, so it is never itself integrated:
\begin{equation}
  \chi(t)=\bigl[1+e^{-k(t-\tcrit)}\bigr]^{-1}.
\end{equation}
At $t=\tcrit$ the switch is exactly one half. The expansion history is
then
\begin{equation}
  E^2(a)=\Omega_{r,0}\,a^{-4}+\Omega_{m,0}\,a^{-3}
  +\Omega_\Lambda\,\chi(t)/\chi_0,
\end{equation}
where $a=1/(1+z)$ is the scale factor and $\chi_0$ is the value of the
switch today, iterated so that $H(a=1)$ equals the sampled $H_0$.

That is the entire phenomenological layer the sampler is allowed to
see. Everything else---distances, supernova magnitudes, BAO ratios,
compressed CMB quantities---is computed from this background.

%------------------------------------------------------------------------------
\section{Materials}
%------------------------------------------------------------------------------

\subsection{Software}

The stack has three layers, plus this report.

\begin{itemize}[leftmargin=1.6em]
  \item \textbf{Layer 1 --- theory.} \texttt{analytical\_solver.py}
        locks the horizon-thermodynamics argument and exports the
        attractor. Inference is forbidden from refitting that attractor.
  \item \textbf{Layer 2 --- Boltzmann.} \texttt{Modified\_CLASS.py}
        is the production Einstein--Boltzmann code. Full CMB power
        spectra live here, not in the sampler.
  \item \textbf{Layer 3 --- inference.} \texttt{Bayesian\_Validationn.py}
        is the production engine: priors, datasets, likelihoods, MCMC,
        nested sampling, diagnostics, plots, and a reproducibility
        snapshot.
  \item \textbf{Layer 3b --- audit.} \texttt{computing\_validation.py}
        checks that the stack is wired correctly and exports this
        manuscript.
\end{itemize}

The default production sampler is a standard-library Markov chain.
Cobaya can be attached as an optional backend; it never owns the
Friedmann equation. Hamiltonian Monte Carlo is declared in the
architecture and raises an error if requested. I did not use it.

A Streamlit booth app can replay a precomputed cartoon of the sampler
and a live $\Delta\chi^2$ scoreboard. That dashboard is a visual aid.
It is not the production inference engine.

\subsection{Data}

Production inference uses public cosmological probes, not synthetic
demo catalogues. Table~\ref{tab:probes} lists the rows the likelihood
knows how to load. Supernova CSVs are read from a \texttt{--data-dir}
folder. Production mode refuses fake posteriors.

\begin{table}[H]
\centering
\caption{MCMC probes used by the joint likelihood. Counts are the
inference sizes stored in the dataset registry, not a claim that every
row was used in every run.}
\label{tab:probes}
\begin{tabular}{llp{0.46\textwidth}}
\toprule
Key & Typical $n$ & What it measures \\
\midrule
Pantheon+      & 1550        & Type~Ia supernova distance moduli $\mu(z)$ \\
DES-SN Y5      & 1635        & Independent SN $\mu(z)$ (cross-check) \\
SH0ES          & 1--40       & Local $H_0$ calibration \\
DESI DR2       & $\sim$10--20 & BAO ratios $D_M/r_d$, $D_H/r_d$, $D_V/r_d$ \\
BOSS DR12      & $\sim$6--12  & Earlier BAO / RSD \\
Planck (compressed) & 3    & Distance priors $(R,\ell_A,\omega_b)$ \\
SPARC / WMAP   & optional    & Rotation curves / multipoles \\
\bottomrule
\end{tabular}
\end{table}

A separate object-count inventory (Type~Ia supernovae, galaxies,
weak-lensing sources, and so on) describes the broader compilation.
Those headline counts are not extra likelihood terms.

%------------------------------------------------------------------------------
\section{Methods}
%------------------------------------------------------------------------------

The work follows a one-way pipeline. Later steps consume earlier
products. Nothing in this list derives Einstein's equation.

\subsection{Lock the physics, then choose what to sample}

I first ran the theory solver and treated its attractor as an input.
The computer is not allowed to refit the late-time Hubble floor
$H_\infty$. I then froze the inference vector to the four numbers in
Section~4. $S_{\max}$ and the early entropy are not sampled.

\subsection{Priors, before any data}

Each parameter gets an independent Gaussian prior
(Table~\ref{tab:priors}). The $H_0$ and $\Omega_\Lambda$ widths come
from algebraic Hubble mixing plus a Planck sound-horizon anchor. The
$k$ and $\tcrit$ widths come from the switch itself: $k$ sets the
transition time $1/k$, and $\tcrit$ is the halfway point.

\begin{table}[H]
\centering
\caption{Table~1 Gaussian priors. These are beliefs \emph{before} the
likelihood, not posterior results.}
\label{tab:priors}
\begin{tabular}{lrrl}
\toprule
Parameter & Centre $\mu$ & Width $\sigma$ & Origin \\
\midrule
$H_0$          & 72.8  & 2.059 & mixing + Planck anchor \\
$\Omega_\Lambda$ & 0.685 & 0.012 & mixing; today's entropy density \\
$k$            & 0.37  & 0.068 & switch speed (Gyr$^{-1}$) \\
$\tcrit$       & 15.8  & 1.408 & halfway time (Gyr) \\
\bottomrule
\end{tabular}
\end{table}

Before MCMC, I drew parameters from this prior alone and generated
ensembles of $H(z)$ and the deceleration history. That prior-predictive
check asks a simple question: does the prior already forbid a late-time
de~Sitter limit? If it did, the later posterior would be untrustworthy
for the wrong reason.

\subsection{Predict the sky from four numbers}

For each trial $\theta$ the background integrator solves for cosmic
time as a function of $\ln a$ with classical Runge--Kutta (RK4),
typically 8000 steps from $a=10^{-3}$ to today. Time is the unknown;
scale factor is the independent variable, so $H(z)$ comes out natively.
A short outer loop updates $\chi_0$ three times so that today's Hubble
rate matches $H_0$ exactly.

Distances are then quadratures of $c/H(z)$. I used composite Simpson's
rule (even number of steps), not a cubic spline, because extra
smoothing would leak into $\chi^2$. From the comoving distance come
the luminosity distance, the distance modulus $\mu$ used by supernovae,
and the BAO combinations $D_M$, $D_H$, and $D_V$.

A theory facade called \texttt{ModifiedCLASS} is the only cosmology
surface the sampler may call. It always runs in the same order:
background, then a small set of Newtonian perturbation modes, then
compressed CMB quantities at last-scattering redshift $z_*=1089.90$,
then a BBKS matter power spectrum, then distances, then linear growth.
Inference is not allowed to poke the raw integrator. Full $C_\ell$
spectra remain Layer~2.

The mandatory sanity check is the $\Lambda$CDM recovery test: when
$k\to 0$,
\begin{equation}
  \max_z \bigl|H_S(z)/H_\Lambda(z)-1\bigr| < 10^{-4}.
\end{equation}
If that test fails, I stop.

\subsection{Score the fit}

Each probe is an independent Gaussian likelihood. If a dataset only
stores error bars, $\chi^2$ is the usual sum of squared residuals. If
it stores a covariance matrix (supernova systematics, BAO), the matrix
form is used. The joint log-likelihood is the sum over probes. A
parameter vector outside the allowed box is assigned
$\ln\pi=-\infty$ and rejected.

\subsection{Sample, then integrate}

I explored the posterior with three optional algorithms:

\begin{itemize}[leftmargin=1.6em]
  \item \textbf{Random-walk Metropolis--Hastings} (default). Propose a
        nearby point; accept it more often when it is more probable.
  \item \textbf{Goodman--Weare stretch moves.} An ensemble of walkers
        stretch toward one another. This is the method behind the
        booth-app walker cartoon.
  \item \textbf{Differential-evolution MCMC.} A walker steps along the
        difference of two others, which helps when parameters are
        correlated.
\end{itemize}

Walkers start at the fiducial values plus a small Gaussian jitter.
Chain lengths, burn-in, walker count, and random seeds are
command-line settings, not hidden constants.

In parallel, nested sampling estimates the evidence $Z=\int L\,\pi\,d\theta$.
A built-in live-point method is always available; dynesty, UltraNest,
or PolyChord are used if they are installed. Evidence is an integral
over all four parameters. It is not a number that should update live
when someone moves a slider on the booth laptop.

\subsection{Refuse to quote a broken chain}

Before any 68\% or 95\% interval is reported, the chain has to pass a
binary \texttt{converged} flag. The checks are standard:

\begin{itemize}[leftmargin=1.6em]
  \item Gelman--Rubin $\widehat R$ (and split-$\widehat R$), default
        fail if $\widehat R>1.1$;
  \item Sokal integrated autocorrelation time and effective sample size,
        default fail if ESS~$<50$ on any parameter;
  \item Geweke's comparison of the first 10\% of a chain with the last
        50\%, used to pick burn-in;
  \item acceptance rate in the window 0.15--0.50, plus stuck-walker
        checks.
\end{itemize}

I do not quote intervals from a chain that fails this flag.

\subsection{Compare models fairly}

At the best-fit point I compute $\chi^2$, AIC, and BIC. Relative to
$\Lambda$CDM, \Lcdms\ has $\Delta k=2$ extra parameters. For Gaussian
likelihoods the penalties are identities:
\begin{equation}
  \Delta\mathrm{AIC}=\Delta\chi^2+4,\qquad
  \Delta\mathrm{BIC}=\Delta\chi^2+2\ln n.
\end{equation}
WAIC, DIC, residual RMS, and a simple leave-one-out style check are
computed automatically. Those scores are useful, but they are not a
substitute for nested-sampling $Z$.

Posterior samples also give means, medians, credible intervals, and
the correlation matrix (including a published pairwise-$r$ table).
$H_0$ is compared with Planck $67.4\pm 0.5$ and SH0ES $73.0\pm 1.0$.
A posterior-predictive check then draws new mock catalogues from the
fitted model and asks whether the real data look typical.

\subsection{Make the run repeatable}

Every production run writes seeds, dataset names, software versions,
a git hash, a 16-character configuration hash, and YAML sidecars.
Anyone with the same data folder and the same command line should be
able to regenerate the chains.

%------------------------------------------------------------------------------
\section{Results}
%------------------------------------------------------------------------------

\subsection{\texorpdfstring{The integrator really is a $\Lambda$CDM extension}{The integrator really is a LambdaCDM extension}}

The Layer-3 audit ran in 0.2~s. Critical checks: 8 passed, 0 failed,
1 skipped (the full unit suite, which is off unless
\texttt{--full} is requested). The important physics check is the
recovery test. At the fiducial background the cosmic age is
$t_0=14.161$~Gyr and today's switch value is $\chi_0=0.3529$. When the
switch is forced off, the maximum fractional Hubble difference versus
flat $\Lambda$CDM is $5.17\times 10^{-6}$, well inside the
$10^{-4}$ budget. The facade distance modulus at $z=0.5$ in that
limit is $\mu=42.165$.

A dummy ``perfect fit'' vector returns $\chi^2=0$, which is the
arithmetic identity a Gaussian likelihood must satisfy. The scope
auditor loaded 11 allowed claims and 10 prohibited slogans, so the
export step itself refuses phrases such as ``MCMC proved the
attractor.''

\begin{table}[H]
\centering
\caption{Live audit of the computing stack. ``Skipped'' means the
optional full Layer-3 suite was not requested.}
\label{tab:audit}
\begin{tabular}{llp{0.52\textwidth}}
\toprule
Status & Check & What it confirmed \\
\midrule
Passed & Import        & \texttt{Bayesian\_Validationn} v2.0.0; 19 phase keys \\
Passed & Layer stack   & Three layers; MCMC samples only the four parameters \\
Passed & Claims        & Allowed / prohibited wording lists are loaded \\
Passed & Parameters    & $\theta=(H_0,\Omega_\Lambda,k,\tcrit)$; Table~1 attached \\
Passed & Recovery      & $\max|\Delta H/H|=5.17\times 10^{-6}$ \\
Passed & Datasets      & 9 MCMC rows; 10 inventory rows \\
Passed & Identities    & Dummy $\chi^2=0$; DES AIC/BIC algebra holds \\
Passed & Facade        & $\mu(z=0.5)=42.165$ in the $\Lambda$CDM limit \\
Skipped & Full suite   & Pass \texttt{--full} to run it \\
\bottomrule
\end{tabular}
\end{table}

\subsection{A slightly better curve is not a better model}

On the published DES comparison with $n=1820$ supernovae, the extra
switch improves the raw fit by
\begin{equation}
  \Delta\chi^2=-2.5
  \quad\Rightarrow\quad
  \Delta\ln L=1.25
  \quad\Rightarrow\quad
  L_S/L_\Lambda \approx 3.49.
\end{equation}
That sounds friendly to \Lcdms\ until the fines are applied. Two extra
parameters give
\[
  \Delta\mathrm{AIC}=+1.5,\qquad \Delta\mathrm{BIC}=+12.5
\]
exactly. Those two numbers are algebra at the published best-fit
point, not a second secret analysis.

The published nested-sampling result is $\Delta\log Z=-6.5$, or a Bayes
factor $B\approx 1.5\times 10^{-3}$ in favour of $\Lambda$CDM. On the
usual Jeffreys wording that is strong to very strong support for the
simpler model. Evidence is the quantity that answers ``which model
should we bet on?''; $\chi^2$ only answers ``which curve hugs these
points more tightly?''

\subsection{What the posterior is allowed to say}

When a production chain passes the convergence flag, the four
parameters receive 68\% and 95\% equal-tail intervals, a covariance
matrix, and an eigensystem that shows the main degeneracies. The
honest scientific sentence is:

\begin{quote}
MCMC constrains $(H_0,\Omega_\Lambda,k,\tcrit)$. Nested sampling
estimates $Z$. Neither calculation derives
$G_{\mu\nu}+\Lambda g_{\mu\nu}=8\pi G T_{\mu\nu}$.
\end{quote}

$H_0$ can be compared with the Planck and SH0ES anchors, but this
pipeline is not a claim that the switch ``solves the Hubble tension.''
That would be one of the prohibited slogans.

%------------------------------------------------------------------------------
\section{Discussion}
%------------------------------------------------------------------------------

\subsection{Why the extra parameters lose}

There is no paradox in ``better $\chi^2$, worse evidence.'' The switch
has two free knobs. On this supernova comparison they buy about two
and a half $\chi^2$ units. AIC already treats that as a wash
($\Delta\mathrm{AIC}=+1.5$ is a weak preference for $\Lambda$CDM).
BIC, which fines harder when $n$ is large, fines this model
noticeably because $n=1820$ is large. Evidence fines even more,
because the prior volume on $(k,\tcrit)$ is not free: most of that
volume does not improve the fit.

That is the scientific posture I want a judge to remember. A more
flexible curve is easy to write down. A more flexible curve that
survives a complexity penalty is rare. On the published numbers,
\Lcdms\ does not survive it.

\subsection{Why the computing still matters}

A negative model-comparison result is still a successful computing
project if the pipeline is honest. The recovery test shows the
integrator is a real $\Lambda$CDM extension, not a broken Hubble
function that happens to have four inputs. The facade forces every
likelihood to see the same theory object. The auditor blocks
overclaims in captions. The reproducibility snapshot means a later
reader can rerun the same command.

Those are the design criteria of the innovation: lock physics, predict
observables, sample only what is allowed, and refuse to quote a
chain that has not mixed.

\subsection{Limitations}

\begin{itemize}[leftmargin=1.6em]
  \item The production facade uses compressed CMB distance priors, not
        a full $C_\ell$ Boltzmann likelihood. Full spectra are Layer~2.
  \item HMC/NUTS is reserved and unimplemented, because it needs
        derivatives of $\ln\pi$ through the theory map.
  \item Independent Gaussian probes are assumed. Shared systematics
        between surveys are not modelled beyond the covariance each
        probe already carries.
  \item Published $\Delta\chi^2$, $\Delta$AIC, $\Delta$BIC, and
        $\Delta\log Z$ are identities or nested-sampling integrals at
        stated settings. They should not be remixed with a different
        supernova catalogue without saying so.
  \item Table~3 object counts describe a compilation. They are not
        thousands of extra $\chi^2$ terms.
\end{itemize}

%------------------------------------------------------------------------------
\section{Conclusions}
%------------------------------------------------------------------------------

I built a three-layer computing stack that tests \Lcdms\ as a
phenomenological extension of flat $\Lambda$CDM. The sampled vector is
only $(H_0,\Omega_\Lambda,k,\tcrit)$. The background integrator
recovers $\Lambda$CDM to $5\times 10^{-6}$ when the extra switch is
turned off. On the published DES comparison the extra switch improves
$\chi^2$ by 2.5, but $\Delta\mathrm{AIC}=+1.5$,
$\Delta\mathrm{BIC}=+12.5$, and $\Delta\log Z=-6.5$ all favour the
simpler model.

The conclusion is therefore two-sided, and both sides are intended:

\begin{enumerate}[leftmargin=1.6em]
  \item \textbf{Physics claim (narrow).} This computing layer does not
        derive de~Sitter from the likelihood. That argument lives in
        the theory file and stands or falls on its own assumptions.
  \item \textbf{Data claim (narrow).} After complexity penalties, these
        probes do not prefer the extra $(k,\tcrit)$ pair.
\end{enumerate}

That is a successful validation of the method, and a negative---but
honest---result for the extra parameters.

%------------------------------------------------------------------------------
\section{What I would do next}
%------------------------------------------------------------------------------

The next useful step is not another slogan. It is a tighter likelihood:
swap the compressed CMB vector for Layer~2 $C_\ell$ spectra, keep the
same four-parameter vector, and rerun nested sampling. If $\Delta\log Z$
stays negative, the extra switch is not helping. If it flips, the
result had better still pass $\widehat R$, ESS, and the recovery test
before anyone quotes it.

%------------------------------------------------------------------------------
\section{Acknowledgements}
%------------------------------------------------------------------------------

Public catalogues from Pantheon+, DES, SH0ES, DESI, BOSS, and Planck
made the likelihood possible. Optional backends were used only when
installed; the standard-library sampler remains the default. Mentors
who read earlier drafts pushed the wording toward what the computer
actually constrained, rather than what would sound impressive.

%------------------------------------------------------------------------------
\section{References}
%------------------------------------------------------------------------------

{\small
\begin{enumerate}[leftmargin=1.8em,itemsep=0.2em]
  \item Planck Collaboration. (2020). Planck 2018 results. VI.
        Cosmological parameters. \textit{Astronomy \& Astrophysics},
        641, A6.
  \item Brout, D., et al. (2022). The Pantheon+ analysis: Cosmological
        constraints. \textit{The Astrophysical Journal}, 938, 110.
  \item DES Collaboration. (2024). The Dark Energy Survey: Cosmology
        results with $\sim 1500$ new high-redshift Type~Ia supernovae
        using the full 5-year dataset. \textit{The Astrophysical Journal
        Letters}, 973, L14.
  \item Riess, A.~G., et al. (2022). A comprehensive measurement of the
        local value of the Hubble constant with 1~km~s$^{-1}$~Mpc$^{-1}$
        uncertainty from the Hubble Space Telescope and the SH0ES team.
        \textit{The Astrophysical Journal Letters}, 934, L7.
  \item DESI Collaboration. (2025). DESI DR2 results. II. Measurements
        of baryon acoustic oscillations and cosmological constraints.
        \textit{arXiv}:2503.14738.
  \item Alam, S., et al. (2017). The clustering of galaxies in the
        completed SDSS-III Baryon Oscillation Spectroscopic Survey.
        \textit{Monthly Notices of the Royal Astronomical Society},
        470, 2617--2652.
  \item Goodman, J., \& Weare, J. (2010). Ensemble samplers with
        affine invariance. \textit{Communications in Applied
        Mathematics and Computational Science}, 5, 65--80.
  \item ter Braak, C.~J.~F. (2006). A Markov chain Monte Carlo version
        of the genetic algorithm Differential Evolution.
        \textit{Statistics and Computing}, 16, 239--249.
  \item Skilling, J. (2006). Nested sampling for general Bayesian
        computation. \textit{Bayesian Analysis}, 1, 833--859.
  \item Gelman, A., \& Rubin, D.~B. (1992). Inference from iterative
        simulation using multiple sequences. \textit{Statistical
        Science}, 7, 457--472.
  \item Akaike, H. (1974). A new look at the statistical model
        identification. \textit{IEEE Transactions on Automatic
        Control}, 19, 716--723.
  \item Schwarz, G. (1978). Estimating the dimension of a model.
        \textit{The Annals of Statistics}, 6, 461--464.
  \item Jeffreys, H. (1961). \textit{Theory of Probability}
        (3rd~ed.). Oxford University Press.
  \item Lewis, A., \& Bridle, S. (2002). Cosmological parameters from
        CMB and other data: A Monte Carlo approach.
        \textit{Physical Review D}, 66, 103511.
\end{enumerate}
}

%------------------------------------------------------------------------------
\appendix
\section{Two-minute judge sheet and how to rerun}
%------------------------------------------------------------------------------

Sentences I will stand behind: I sample four numbers, not the attractor.
Turning the switch off recovers $\Lambda$CDM to $5.17\times 10^{-6}$.
On DES, $\chi^2$ improves by 2.5, but AIC, BIC, and evidence still favour
$\Lambda$CDM. $\Delta\log Z=-6.5$ is an integral, not a slider. The
Streamlit app is a cartoon; the sampler is
\texttt{Bayesian\_Validationn.py}. Do not write ``MCMC proved the
attractor.''

All MCMC sizes are command-line flags, for example:
\texttt{python Bayesian\_Validationn.py --prod-steps 2000 --burn 400
--chains 4 --data-dir PATH --require-sn}, and
\texttt{python computing\_validation.py} for the audit that exported
this report.

\end{document}
